\textbf{COLUNI 2023 -- Questão 06}\par\medskip Analise o trecho:

\par

“Acenou \underline{nervosamente} para o motorista para que desviasse, e aí nova surpresa: o homem também lhe acenava, com o mesmo propósito.” Na contramão da história. SCLIAR, Moacyr. \textbf{Histórias que os jornais não contam}. 3. ed. Porto Alegre: L\&PM Editores, 2018. p. 9. (Adaptado).

\par

Verifique as alternativas que expressam a mesma circunstância do termo em destaque no trecho acima:

\par

I.     “Uma alteração tão significativa deveria ter sido \underline{previamente} divulgada; e teria sido necessário colocar avisos na rodovia.” (SCLIAR, Moacyr. O rádio apaixonado. In: SCLIAR, Moacyr. \textbf{Histórias que os jornais não contam}. 3. ed. Porto Alegre: L\&M Editores, 2018. p.9).\par
II. “Enfrentara \underline{corajosamente} o julgamento, sem mudar de opinião. E ele não mudaria de pista.” (SCLIAR, Moacyr. O rádio apaixonado. In: SCLIAR, Moacyr. \textbf{Histórias que os jornais não contam}. 3. ed. Porto Alegre: L\&M Editores, 2018. p.10).\par
III. “Um dia, porém, a VERDADE apareceria naquela estrada. Avançando \underline{celeremente}, e na mesma mão em que ele estava.” (SCLIAR, Moacyr. O rádio apaixonado. In: SCLIAR, Moacyr. \textbf{Histórias que os jornais não contam}. 3. ed. Porto Alegre: L\&M Editores, 2018. p.9).

\par

Está CORRETO o que se afirma em:\begin{enumerate}
\item I e II, apenas.
\item II e III, apenas.
\item I, apenas.
\item II, apenas.\end{enumerate}
